\documentclass[11pt]{article}
\usepackage[margin=25mm]{geometry}
\usepackage{fontspec}
\setmainfont{Latin Modern Roman}
\usepackage{amsmath,amssymb,mathtools}
\usepackage{unicode-math}
\setmathfont{Latin Modern Math}
\usepackage{xurl,hyperref}
\hypersetup{colorlinks=true,urlcolor=blue}
\setcounter{secnumdepth}{0}
\setlength{\parindent}{0pt}
\setlength{\parskip}{5pt}
\setlength{\emergencystretch}{3em}
\newcommand{\sref}[2]{\hyperlink{source:#1:#2}{[S#2]}}
\newcommand{\cataloguescope}[1]{\paragraph{Recorded target.}#1}
\newcommand{\R}{\mathbb R}
\newcommand{\C}{\mathbb C}
\newcommand{\N}{\mathbb N}
\newcommand{\Z}{\mathbb Z}
\newcommand{\Q}{\mathbb Q}
\newcommand{\F}{\mathbb F}
\newcommand{\eps}{\varepsilon}
\begin{document}
% UnsolvedMath ID 30003751; source catalogue code OWR-15998-001
\setcounter{section}{30003750}
\section{Regularized Density-to-Potential Maps in Density Functional Theory}\label{30003751}
{\small UnsolvedMath ID 30003751 · Mathematical Physics}
\subsection{Definitions and mathematical statement}

\paragraph{Shared notation.}$\mathbb N=\{1,2,\ldots\}$; the standard number fields have their usual meanings. All additional conventions are stated locally.


\paragraph{Definitions and conventions.}
For prescribed particle density $\rho$, the Levy--Lieb functional is a constrained infimum of kinetic plus interaction energy over quantum states with that density. A density-to-potential map seeks an external potential realizing the density in the relevant variational problem, modulo an additive constant. Given a Banach or Hilbert density space $X$ and a norm, the Moreau--Yosida envelope has the form $F_\varepsilon(\rho)=\inf_{\sigma\in X}\{F(\sigma)+(2\varepsilon)^{-1}\|\rho-\sigma\|_X^2\}$; differentiability depends on the space and convex formulation.

\paragraph{Target retained from the canonical source catalogue.}
\paragraph{Catalogue wording retained for traceability.}
Establish rigorous conditions for existence of the density-to-potential map in density functional theory, and use functional-analytic Moreau–Yosida regularization to address the nondifferentiability of the constrained-search Levy–Lieb functional.

\paragraph{Source scope and qualification.}
The chosen pure-state or ensemble functional, density space and potential dual space are essential. A regularized density-to-potential map is not silently identified with an unregularized map for every representable density.
\subsection{Short English statement}
Specify when densities determine potentials and use rigorous regularization to handle nondifferentiable density functionals.
\subsection{Research attempt}
\paragraph{Model, date and investigated scope.}
OpenAI GPT-6 Astra (configured model \texttt{gpt-6-astra}), ultra reasoning; generated 11 October 2026. Computed a Moreau envelope and showed unique regularized gradients do not establish unique unregularized density-to-potential maps.

\paragraph{Sources and literature check.}
Read the supplied definition, catalogue wording, source qualifications and reference entries for 30003751. The historical source/context and inherited status are retained; no fresh literature or novelty certification is claimed. No external theorem is used without proof in this attempt; explicit model assumptions are distinguished from facts about the source model.

\paragraph{A complete scalar regularization test.} To test the inference from a regularized derivative to a unique physical potential, use the convex function \(F(x)=|x|\) on \(\mathbb R\). Minimizing \(|y|+(x-y)^2/(2\varepsilon)\), \(\varepsilon>0\), separately on the positive and negative half-lines gives
\[
 y_\varepsilon(x)=\operatorname{sgn}(x)(|x|-\varepsilon)_+,
 \qquad
 F_\varepsilon(x)=\begin{cases}x^2/(2\varepsilon),&|x|\le\varepsilon,\\|x|-\varepsilon/2,&|x|\ge\varepsilon.\end{cases}
\]
The two values and their first derivatives agree at \(\pm\varepsilon\); hence \(F_\varepsilon'\) is the clipped function \(x/\varepsilon\) in the central interval and \(\operatorname{sgn}(x)\) outside. It is unique everywhere and Lipschitz with constant \(1/\varepsilon\).
\paragraph{What this does not identify.} At \(x=0\), the unregularized variational problem \(\min_x(|x|+vx)\) has zero as a minimizer for every \(v\in[-1,1]\), since \(|x|+vx\ge(1-|v|)|x|\). The regularized derivative chooses zero there but does not prove uniqueness of an unregularized potential. At the endpoints \(v=\pm1\), even the minimizing density is nonunique along a half-line.
\paragraph{Source comparison.} This known convex-analysis example isolates the precise logical gap in an inversion route. It is not a Levy--Lieb functional and proves no physical representability claim. The source requires a specified pure-state or ensemble density space, its potential dual and hypotheses transferring regularization to physical densities modulo constants. Proving that transfer, especially as \(\varepsilon\downarrow0\), remains the first gap; the differentiability calculation alone cannot supply it.

\paragraph{Verification and adversarial audit.}
Checked endpoint differentiability, the Lipschitz constant, and all potentials in the scalar subgradient interval. This is a GPT-6 Astra Ultra self-audit, not independent human verification. Prompt A drafting followed by prompt B scope/dependency/boundary self-audit. Checked endpoint differentiability, the Lipschitz constant, and all potentials in the scalar subgradient interval. The remaining gap is: Establish representability and the physical dual-space limit for the chosen Levy--Lieb formulation. Unresolved status and zero original-target completion are retained.

Coordinator checked the soft-threshold minimizer, differentiability at both joins, and nonuniqueness of the original potential at zero. This additional check was model-assisted, not independent human verification.

\paragraph{Tools, process and contributions.}
Python 3 and local file-reading/writing tools for manuscript serialization and structural checks; mathematical calculations were symbolic reasoning, with no numerical experiment or theorem prover. Contribution prompt A was followed by prompt B. GPT-6 Astra with ultra reasoning drafted the research and performed its stated self-audit; any additional coordinator contribution is identified in the verification and audit. A GPT-6-based Codex coordinator assembled the manuscript and checked repository integration; its exact variant and reasoning configuration were not exposed. No human mathematical contribution, independent human proof check or formal proof certification is supplied. Complete derivations and the consulted source references are included here; private reasoning transcripts are not published.

\paragraph{Final status and first remaining gap.}
\textbf{UNRESOLVED}. Establish representability and the physical dual-space limit for the chosen Levy--Lieb formulation.

\section{Completion Estimate}
\noindent\textbf{COMPLETION: 0\%}
The proved calculation is an elementary auxiliary result or a test of a restricted strategy; no new portion of the full source target is established.

\subsection{Sources}
\begin{itemize}
\item[\hypertarget{source:30003751:1}{S1}] ULAM.AI and the UnsolvedMath Contributors, UnsolvedMath, record 30003751 (OWR-15998-001). Catalogue curation under CC BY 4.0; definitions and English exposition added locally.
\url{https://huggingface.co/datasets/ulamai/UnsolvedMath}.
\item[\hypertarget{source:30003751:2}{S2}] Mathematical Methods in Quantum Chemistry (2018), Oberwolfach Reports, DOI 10.4171/owr/2018/13, pp. 631--721. Primary report accessed 2026-10-11; relevant PDF pages 2.
\url{https://doi.org/10.4171/owr/2018/13}.
\item[\hypertarget{source:30003751:3}{S3}] CDFT Part I (catalogue-listed context; not independently consulted in this preparation).
\url{https://arxiv.org/abs/2211.16627}.
\item[\hypertarget{source:30003751:4}{S4}] CDFT Part II (catalogue-listed context; not independently consulted in this preparation).
\url{https://arxiv.org/abs/2303.01357}.
\item[\hypertarget{source:30003751:5}{S5}] Periodic regularized inversion (catalogue-listed context; not independently consulted in this preparation).
\url{https://journals.aps.org/prb/abstract/10.1103/PhysRevB.111.205143}.
\end{itemize}
\label{end:30003751}
\end{document}
